\subsection{Complements on Polish spaces}

Lemma 1. --- Let X be a topological space and let A be a subset of X. For A to be meager, it is necessary and sufficient that there exists an open covering \((\mathrm{U}_{i})_{i\in\mathrm{I}}\) of X such that \(A\cap U_{i}\) be meager in \(U_{i}\) for all \(i\in I\) .

Let \(\mathcal{O}\) be the set of open subsets U of X such that \(A \cap U\) is meager. The family of subsets of \(\mathcal{O}\) consisting of pairwise disjoint open sets, ordered by inclusion, is inductive. Let \(\mathfrak{U}\) be a maximal element; such an element exists by E, III, p. 20, th. 2.

Let O be the union of the open sets of X belonging to \(\mathfrak{U}\). For every open set \(U \in \mathfrak{U}\) , let \((B_{U,n})\) be a sequence of nowhere dense subsets of U whose union is \(A \cap U\) . For every integer n, let \(B_n = \bigcup_{U \in \mathfrak{U}} B_{U,n}\) .

As U ranges over U, it is nowhere dense relative to O (TG, IX, p. 52, prop. 1). By TG, IX, p. 53, prop. 2, \(B_{n}\) is a nowhere dense subset of X, since O is open in X. Consequently, the set A ∩ O, equal to the union of the \(B_{n}\) is a meager subset of X.

Let F be the complement of O. It is a closed subset of X; let us show that it has empty interior. Otherwise, let x be an interior point of F. By hypothesis, there exists an open neighborhood V of x, which may be assumed to be contained in F, such that A ∩ V is meager. Then V is disjoint from the open sets of X belonging to U, and U ∪ \{V\} is a set of pairwise disjoint open sets belonging to O, which contradicts the maximality of U. The relation

\[
\mathrm{A} = (\mathrm{A} \cap \mathrm{F}) \cup (\mathrm{A} \cap \mathrm{O})
\]

then implies that A is meager, which is what was to be proved.

Let X be a topological space.

Recall (TG, IX, p. 69) that a subset A of X is said to be approachable if there exists an open set U of X such that \(U \cap \complement A\) and \(A \cap \complement U\) are meager in X. The collection of approachable subsets of X is a \(\sigma\)-algebra that contains the Borel \(\sigma\)-algebra (TG, IX, p. 69, Lemma 8 and its proof). A meager subset is approachable.

For any subset A of X, let D(A) be the set of points \(x \in X\) such that for every neighborhood V of x, A∩V is not meager. We also set \(\mathrm{D}^{*}(\mathrm{A}) = \mathrm{A} \cup \mathrm{D}(\mathrm{A})\) .

Lemma 2. --- Let X be a topological space and let A be a subset of X.

\begin{enumerate}
\def\labelenumi{\alph{enumi})}
\item
  The set D(A) is closed in X; its complement is the largest open set U in X such that A ∩ U is a meager set.
\item
  For A to be meager, it is necessary and sufficient that D(A) be empty.
\item
  The set \(\mathrm{D}^{*}(\mathrm{A})\) is approachable.
\item
  For every approachable set B of X containing A, \(D^{*}(A) \cap \complement B\) is meager.
\end{enumerate}

Let U be the union of the open sets V of X such that \(A \cap V\) is meager. For a point x to belong to U, it is necessary and sufficient that it have a neighborhood V such that \(A \cap V\) is meager. Thus \(U = \complement D(A)\) , and D(A) is therefore a closed subset of X. By construction, the subspace U has a covering by open sets whose intersection with A is meager. By Lemma 1, applied to the topological space U and the subset \(A \cap U\) , \(A \cap U\) is meager in U, hence also in X. This proves assertion a), since, by construction, every open set V of X such that \(A \cap V\) is meager is contained in U.

Assertion b) follows immediately from this.

\begin{enumerate}
\def\labelenumi{\alph{enumi})}
\setcounter{enumi}{2}
\item
  We have \(D^{*}(A) = A \cup D(A) = (A \cap U) \cup D(A)\) . The set \(D(A)\) is closed, hence approachable (IV, p. 361); the meager part \(A \cap U\) is also approachable. Consequently, \(D^{*}(A)\) is approachable (loc. cit.).
\item
  Let B be an approachable subset of X containing A. Its complement \(\complement B\) is then an approachable subset of X (loc. cit.), and there therefore exists an open set V of X such that \(V \cap B\) and \(\complement V \cap \complement B\) are meager. Since \(A \subset B\) , \(V \cap A\) is still meager, whence \(V \subset U\) . Since \(\complement U \cap \complement B\) is contained in \(\complement V \cap \complement B\) , it is also a meager subset of X. Finally, the inclusions
\end{enumerate}

\[
\mathrm{D} ^ {*} (\mathrm{A}) \cap \complement \mathrm{B} = ((\mathrm{A} \cap \mathrm{U}) \cap \complement \mathrm{B}) \cup (\complement \mathrm{U} \cap \complement \mathrm{B}) \subset (\mathrm{A} \cap \mathrm{U}) \cup (\complement \mathrm{U} \cap \complement \mathrm{B})
\]

prove that \(D^{*}(A) \cap \complement B\) is a meager subset of X.

Remark 1. --- Let G be a topological group and let A be a meager subset of G. Suppose that A contains a nonempty open set U of G, and let y be a point of U. Then U is meager, and for every \(x \in G\) , \(xy^{-1}U\) is a neighborhood of x in G. Consequently, every point of G has a meager neighborhood, so G is meager (IV, p. 360, lemma 1).

Conversely, if G is not meager, every meager subset has empty interior and G is a Baire space. Since a Baire space is not meager, this proves the following assertion: For a topological group to be a Baire space, it is necessary and sufficient that it be not meager.

PROPOSITION 7. --- Let X be a Hausdorff space. Every Souslin subspace (TG, IX, p. 59, def. 2) of X is approachable.

Let S be a Souslin subspace of X. By definition, there exists a complete metric space with countable base P and a continuous surjective map g from P onto S. By TG, IX, p. 64, lemma 3, there exists a sieve \(\mathrm{C} = (\mathrm{C}_{n}, p_{n}, \varphi_{n})_{n \in \mathbf{N}}\) of the metric space P.

For every integer \(n\) and every \(c \in \mathrm{C}_n\) , denote by \(\mathrm{F}_n(c)\) the image of the set \(\varphi_n(c)\) under \(g\) ; let us also set \(\mathrm{F}_n^*(c) = \mathrm{D}^*(\mathrm{F}_n(c))\) and

\[
\mathrm{G} _ {n} (c) = \mathrm{F} _ {n} ^ {*} (c) \cap \mathbb {C} \left(\bigcup_ {c ^ {\prime} \in p _ {n} ^ {- 1} (c)} \mathrm{F} _ {n + 1} ^ {*} (c ^ {\prime})\right).\tag{1}
\]

For every \(c \in C_{n}\) , \(\mathrm{F}_{n}^{*}(c)\) is approachable (IV, p. 361, lemma 2, c)). Since \(C_{n+1}\) is countable, the union of the approachable subsets \(\mathrm{F}_{n+1}^{*}(c')\) , for \(c'\) ranging over \(p_{n}^{-1}(c)\) , is again an approachable subset of X. It contains the union of the \(\mathrm{F}_{n+1}(c')\) , which equals \(\mathrm{F}_{n}(c)\) . Consequently (loc. cit., d)), \(\mathrm{G}_{n}(c)\) is a meager subset of X. The union G of the sets \(\mathrm{G}_{n}(c)\) , for \(n \in \mathbf{N}\) and \(c \in C_{n}\) , is therefore a meager subset of X.

Let \(c_0 \in \mathrm{C}_0\) and let \(x\) be an element of \(\mathrm{F}_0^*(c_0) \cap \complement \mathrm{G}\) ; let us prove that \(x \in \mathrm{F}_0(c_0)\) . Since \(x \notin \mathrm{G}_0(c_0)\) and \(x \in \mathrm{F}_0^*(c_0)\) , there exists \(c_1 \in p_0^{-1}(c_0)\) such that \(x \in \mathrm{F}_1^*(c_1)\) by relation (1). By induction, there exists an element \(c = (c_n)_n\) of \(\prod_n C_n\) such that, for every integer \(n \in \mathbf{N}\) , \(x \in \mathrm{F}_n^*(c_n)\) and \(p_n(c_{n+1}) = c_n\) .

The subsets \(\varphi_n(c_n)\) form a base of a Cauchy filter in P; this filter converges to a point \(p\) of P, since P is complete. The image of this filter under \(g\) is a filter F on X, with base the set of the \(\mathrm{F}_n(c_n)\) , for \(n\in \mathbf{N}\) , which converges to \(g(p)\) . Since \(\mathrm{F}_n^* (c_n)\) is contained in \(\overline{\mathrm{F}_n(c_n)}\) and \(x\) belongs to each of the \(\mathrm{F}_n(c_n)\) , \(x\) is an adherent point of F, hence \(x = g(p)\) since the space X is Hausdorff (TG, I, p. 52, prop. 1). The point \(p\) belongs to \(\overline{\varphi_1(c_1)}\) , hence to \(\varphi_0(c_0)\) . Consequently, \(x\in \mathrm{F}_0(c_0)\) .

Consequently, the set \(\mathrm{F}_{0}^{*}(c_{0}) - \mathrm{F}_{0}(c_{0})\) is contained in G. It is thus a meager subset of X, hence also an approachable subset. Since \(\mathrm{F}_{0}(c_{0}) = \mathrm{F}_{0}^{*}(c_{0}) - (\mathrm{F}_{0}^{*}(c_{0}) - \mathrm{F}_{0}(c_{0}))\) and \(\mathrm{F}_{0}^{*}(c_{0})\) is approachable, \(\mathrm{F}_{0}(c_{0})\) is approachable, because the approachable subsets of X form a \(\sigma\)-algebra. Since \(C_{0}\) is countable, the set S which is the union of the \(\mathrm{F}_{0}(c_{0})\) , for \(c_{0} \in C_{0}\) , is still approachable.

